% ============================================================================
% CHAPTER 1G: AEROSPACE SYSTEMS
% ============================================================================

\chapter{Aerospace Systems}
\label{ch:aerospace}

Aerospace systems present unique control challenges: operation in three dimensions, coupling between axes, varying dynamics with flight conditions, and stringent safety requirements.

% ============================================================================
\section{Aircraft Longitudinal Dynamics}
\label{sec:longitudinal}
% ============================================================================

\input{figures/fig_ch13_aircraft_longitudinal.tex}

\subsection{State Variables}

\begin{itemize}
    \item $u$ = forward velocity perturbation (m/s)
    \item $w$ = vertical velocity perturbation (m/s)
    \item $q$ = pitch rate (rad/s)
    \item $\theta$ = pitch angle (rad)
\end{itemize}

\subsection{Linearized Equations of Motion}

\begin{equation}
\begin{bmatrix} \dot{u} \\ \dot{w} \\ \dot{q} \\ \dot{\theta} \end{bmatrix} =
\begin{bmatrix}
X_u & X_w & 0 & -g \\
Z_u & Z_w & U_0 & 0 \\
M_u & M_w & M_q & 0 \\
0 & 0 & 1 & 0
\end{bmatrix}
\begin{bmatrix} u \\ w \\ q \\ \theta \end{bmatrix} +
\begin{bmatrix} X_{\delta_e} \\ Z_{\delta_e} \\ M_{\delta_e} \\ 0 \end{bmatrix} \delta_e
\end{equation}

where $\delta_e$ is elevator deflection and $X_u$, $Z_w$, etc. are stability derivatives.

\subsection{Characteristic Modes}

The longitudinal dynamics typically have two modes:

\subsubsection{Short Period Mode}
\begin{itemize}
    \item Fast oscillation (1-3 seconds period)
    \item Primarily $\alpha$ and $q$ motion
    \item Well-damped ($\zeta \approx 0.3-0.7$)
\end{itemize}

Approximate transfer function:
\begin{equation}
\frac{\alpha(s)}{\delta_e(s)} \approx \frac{Z_{\delta_e}/U_0}{s^2 - (M_q + M_{\dot{\alpha}} + Z_w/U_0)s + (Z_w M_q/U_0 - M_w)}
\end{equation}

\subsubsection{Phugoid Mode}
\begin{itemize}
    \item Slow oscillation (30-100 seconds period)
    \item Exchange between kinetic and potential energy
    \item Lightly damped ($\zeta \approx 0.04-0.1$)
\end{itemize}

Approximate frequency and damping:
\begin{align}
\omega_p &\approx \sqrt{2}\frac{g}{U_0} \\
\zeta_p &\approx \frac{1}{\sqrt{2}}\frac{C_D}{C_L}
\end{align}

\begin{examplebox}[title=Example 1.14: Boeing 747 Longitudinal Modes]
At cruise ($U_0 = 250$ m/s, altitude 10 km):

Short period:
\begin{itemize}
    \item $\omega_{sp} = 0.9$ rad/s, $\zeta_{sp} = 0.35$
    \item Period: $T_{sp} = 2\pi/\omega_{sp} = 7$ s
\end{itemize}

Phugoid:
\begin{itemize}
    \item $\omega_p = 0.06$ rad/s, $\zeta_p = 0.05$
    \item Period: $T_p = 2\pi/\omega_p = 105$ s
\end{itemize}
\end{examplebox}

% ============================================================================
\section{Aircraft Body Axes and Rotational Dynamics}
\label{sec:body_axes}
% ============================================================================

\input{figures/fig_ch13_aircraft_body_axes.tex}

% ============================================================================
\section{Aircraft Lateral-Directional Dynamics}
\label{sec:lateral}
% ============================================================================

\subsection{State Variables}

\begin{itemize}
    \item $\beta$ = sideslip angle (rad)
    \item $p$ = roll rate (rad/s)
    \item $r$ = yaw rate (rad/s)
    \item $\phi$ = bank angle (rad)
\end{itemize}

\subsection{Characteristic Modes}

\subsubsection{Dutch Roll}
\begin{itemize}
    \item Coupled yaw-roll oscillation
    \item Period: 3-10 seconds
    \item Damping varies with aircraft design
\end{itemize}

\subsubsection{Roll Mode}
\begin{itemize}
    \item Pure rolling motion
    \item First-order, typically fast (time constant 0.5-2 s)
\end{itemize}

\subsubsection{Spiral Mode}
\begin{itemize}
    \item Slow divergence or convergence
    \item Often slightly unstable (acceptable if slow)
\end{itemize}

% ============================================================================
\section{Rocket Dynamics}
\label{sec:rocket}
% ============================================================================

\input{figures/fig_ch13_rocket.tex}

\subsection{Pitch Dynamics (Simplified)}

For a rocket with thrust vector control:
\begin{equation}
I_{yy}\ddot{\theta} = T \cdot \ell_T \cdot \delta
\end{equation}

where:
\begin{itemize}
    \item $I_{yy}$ = pitch moment of inertia
    \item $T$ = thrust
    \item $\ell_T$ = distance from CG to gimbal point
    \item $\delta$ = gimbal angle
\end{itemize}

Transfer function:
\begin{equation}
\frac{\Theta(s)}{\delta(s)} = \frac{T\ell_T/I_{yy}}{s^2} = \frac{K_r}{s^2}
\end{equation}

This is a double integrator---marginally stable and requires feedback for control.

\subsection{Aerodynamic Effects}

With aerodynamic forces, additional terms appear:
\begin{equation}
I_{yy}\ddot{\theta} = T\ell_T\delta + M_\alpha\alpha + M_q q
\end{equation}

For a statically unstable rocket ($M_\alpha > 0$), the system has a pole in the right half-plane.

% ============================================================================
\section{Satellite Attitude Dynamics}
\label{sec:satellite}
% ============================================================================

\input{figures/fig_ch13_satellite.tex}

\subsection{Euler's Equations}

\begin{align}
I_{xx}\dot{\omega}_x + (I_{zz} - I_{yy})\omega_y\omega_z &= \tau_x \\
I_{yy}\dot{\omega}_y + (I_{xx} - I_{zz})\omega_z\omega_x &= \tau_y \\
I_{zz}\dot{\omega}_z + (I_{yy} - I_{xx})\omega_x\omega_y &= \tau_z
\end{align}

\subsection{Linearized Single-Axis Model}

For small angles about one axis:
\begin{equation}
I\ddot{\theta} = \tau
\end{equation}

Transfer function:
\begin{equation}
\frac{\Theta(s)}{\tau(s)} = \frac{1}{Is^2}
\end{equation}

\subsection{Reaction Wheel Dynamics}

A reaction wheel applies torque by changing its angular momentum:
\begin{equation}
\tau_{\text{satellite}} = -I_w\dot{\omega}_w
\end{equation}

Combined system:
\begin{equation}
I_s\ddot{\theta} = -I_w\dot{\omega}_w = I_w\ddot{\theta}_w
\end{equation}

% ============================================================================
\section{Quadrotor Dynamics}
\label{sec:quadrotor}
% ============================================================================

\input{figures/fig_ch13_quadrotor_fbd.tex}

\input{figures/fig_ch13_quadrotor.tex}

\subsection{Control Inputs}

\begin{align}
\text{Total thrust:} \quad T &= F_1 + F_2 + F_3 + F_4 \\
\text{Roll torque:} \quad \tau_\phi &= \ell(F_1 - F_2) \\
\text{Pitch torque:} \quad \tau_\theta &= \ell(F_3 - F_4) \\
\text{Yaw torque:} \quad \tau_\psi &= k_\tau(F_1 + F_2 - F_3 - F_4)
\end{align}

\subsection{Equations of Motion}

Translational (in body frame):
\begin{align}
m\ddot{x} &= -mg\sin\theta \\
m\ddot{y} &= mg\cos\theta\sin\phi \\
m\ddot{z} &= mg\cos\theta\cos\phi - T
\end{align}

Rotational:
\begin{align}
I_{xx}\ddot{\phi} &= \tau_\phi \\
I_{yy}\ddot{\theta} &= \tau_\theta \\
I_{zz}\ddot{\psi} &= \tau_\psi
\end{align}

\subsection{Linearized Model (Hover)}

Near hover ($\phi, \theta, \psi \approx 0$):
\begin{equation}
\ddot{\phi} = \frac{\ell}{I_{xx}}(F_1 - F_2), \quad
\ddot{\theta} = \frac{\ell}{I_{yy}}(F_3 - F_4), \quad
\ddot{z} = g - \frac{T}{m}
\end{equation}

% ============================================================================
\section{Spacecraft Orbital Mechanics}
\label{sec:orbital}
% ============================================================================

\subsection{Hill-Clohessy-Wiltshire Equations}

Relative motion near a circular reference orbit:
\begin{align}
\ddot{x} - 2n\dot{y} - 3n^2 x &= a_x \\
\ddot{y} + 2n\dot{x} &= a_y \\
\ddot{z} + n^2 z &= a_z
\end{align}

where $n = \sqrt{\mu/a^3}$ is the orbital rate.

The out-of-plane motion ($z$) is decoupled and is a simple harmonic oscillator.

The in-plane motion ($x$, $y$) is coupled, with the solution:
\begin{equation}
x(t) = (4 - 3\cos nt)x_0 + \frac{\sin nt}{n}\dot{x}_0 + \frac{2(1-\cos nt)}{n}\dot{y}_0
\end{equation}

% ============================================================================
